% Laget automatisk av elevark/lag_elevark.py fra sim/loddrett-sirkel.html. Endringer her blir overskrevet.
\documentclass[11pt]{article}
\usepackage{elevark}
\begin{document}
\elevarkhode{Fysikk 2 · Sirkelbevegelse}{Vertikal sirkelbevegelse}{https://fysikklab.no/sim/loddrett-sirkel.html}{Simuleringen}

En ball svinges rundt i en vertikal sirkel, eller en vogn kjører gjennom en loop i en berg-og-dal-bane. Farten er størst i bunnen og minst på toppen. Er farten for liten, blir snora slakk eller vogna slipper banen. Hvor stor fart trengs i bunnen for å komme helt rundt?

\section*{Oppgaver}

Start med ball i snor, $r = 1{,}0\,\mathrm{m}$, 7,5~m/s i bunnen og 0,50~kg.

\begin{oppgaveliste}
\oppgave \fase{gjett}{Gjett}Hvor er snorkraften størst, og hvor er den minst? Se på grafen.
\svar
\oppgave \fase{test}{Regn}Regn ut snorkraften i bunnen. Bruk energibevaring til å finne farten på toppen, og regn ut snorkraften der. Sjekk.
\svar
\oppgave \fase{forklar}{Forklar}Vis at ballen akkurat kommer rundt hvis farten på toppen er $v = \sqrt{gr}$. Bruk energibevaring til å vise at farten i bunnen da må være $\sqrt{5gr}$.
\svar
\oppgave \fase{test}{Undersøk}Sett farten i bunnen til 6,0~m/s. Hvor blir snora slakk? Hva gjør ballen etterpå?
\svar
\oppgave \fase{test}{Undersøk}Sett farten til 4,0~m/s. Hvorfor blir ikke snora slakk nå? Hvor høyt kommer ballen?
\svar
\oppgave \fase{forklar}{Forklar}Bytt til vogn i loop. Hva er forskjellen og likheten mellom snorkraften og normalkraften i de to situasjonene? Hvorfor kjenner du deg tung i bunnen av en loop og lett på toppen?
\svar
\oppgave \fase{utfordring}{Utfordring}En berg-og-dal-bane har en loop med radius 8,0~m. Hvor stor fart må vogna minst ha i bunnen? Hvor mange $g$ kjenner passasjerene i bunnen da?
\svar
\end{oppgaveliste}

\end{document}
